\scanpage{58}
\mspage{VI}

\deleted{On est} Récurrence sur $\dim X$, \inserted{trivial si elle est $\unsure{\leq|<}0$.} \deleted{\unsure{Utilisant}} \deleted{On \unsure{a}} Lemme 3, on est ramené au cas où $X$ est lisse et \deleted{\unsure{admet} un morphisme lisse surjectif} \inserted{connexe. Si $\dim X=1$,} c'est facile (réduction au cas complet par lemme de Abhyankar) \deleted{\unsure{cf}} c'est dû à \unsure{a} Serre-Lang. \deleted{\unsure{Dim supérieure}}

Si $\dim X\geq 2$, on est ramené \deleted{\unsure{au cas}} \inserted{par} lemme 3 au cas où il y a une fibration \inserted{lisse} $X\longrightarrow Y$, $\dim Y<\dim X$, \deleted{et} si on veut à fibres \unsure{géom.|géom} connexes non vides, de dim 1. On applique alors le lemme 1 et l'hyp. de récurrence.
\margin{?}
\uline{Prop} Soit $k$ un corps. On suppose que pour tte \deleted{\unsure{$X$}} courbe \inserted{alg.} lisse $X$ sur $k$, le \unsure{nb} des rev. étales de $X$ de degré donné soit fini. Alors la $=$ chose est vraie pour \uline{tt} $X$ de t. f. sur $k$.
% archive sheet number: 57
